\documentclass[10pt,a4paper]{amsart}
\usepackage{amsmath,amssymb,amsthm,graphicx}
\usepackage{verbatim}
\newenvironment{reference}{\comment}{\endcomment}
%\newenvironment{reference}{}{}
\newenvironment{slogan}{\comment}{\endcomment}
\newenvironment{history}{\comment}{\endcomment}

% For commutative diagrams we use Xy-pic
\usepackage[all]{xy}

% We use 2cell for 2-commutative diagrams.
\xyoption{2cell}
% Theorem environments.
%
\theoremstyle{plain}
\newtheorem{theorem}[subsection]{Theorem}
\newtheorem{proposition}[subsection]{Proposition}
\newtheorem{lemma}[subsection]{Lemma}

\theoremstyle{definition}
\newtheorem{definition}[subsection]{Definition}
\newtheorem{example}[subsection]{Example}
\newtheorem{exercise}[subsection]{Exercise}
\newtheorem{situation}[subsection]{Situation}

\theoremstyle{remark}
\newtheorem{remark}[subsection]{Remark}
\newtheorem{remarks}[subsection]{Remarks}

\numberwithin{equation}{subsection}

% Macros
%
\def\lim{\mathop{\mathrm{lim}}\nolimits}
\def\colim{\mathop{\mathrm{colim}}\nolimits}
\def\Spec{\mathop{\mathrm{Spec}}}
\def\Hom{\mathop{\mathrm{Hom}}\nolimits}
\def\Ext{\mathop{\mathrm{Ext}}\nolimits}
\def\SheafHom{\mathop{\mathcal{H}\!\mathit{om}}\nolimits}
\def\SheafExt{\mathop{\mathcal{E}\!\mathit{xt}}\nolimits}
\def\Sch{\mathit{Sch}}
\def\Mor{\mathop{\mathrm{Mor}}\nolimits}
\def\Ob{\mathop{\mathrm{Ob}}\nolimits}
\def\Sh{\mathop{\mathit{Sh}}\nolimits}
\def\NL{\mathop{N\!L}\nolimits}
\def\CH{\mathop{\mathrm{CH}}\nolimits}
\def\proetale{{pro\text{-}\acute{e}tale}}
\def\etale{{\acute{e}tale}}
\def\QCoh{\mathit{QCoh}}
\def\Ker{\mathop{\mathrm{Ker}}}
\def\Im{\mathop{\mathrm{Im}}}
\def\Coker{\mathop{\mathrm{Coker}}}
\def\Coim{\mathop{\mathrm{Coim}}}

% Boxtimes
%
\DeclareMathSymbol{\boxtimes}{\mathbin}{AMSa}{"02}

%
% Macros for moduli stacks/spaces
%
\def\QCohstack{\mathcal{QC}\!\mathit{oh}}
\def\Cohstack{\mathcal{C}\!\mathit{oh}}
\def\Spacesstack{\mathcal{S}\!\mathit{paces}}
\def\Quotfunctor{\mathrm{Quot}}
\def\Hilbfunctor{\mathrm{Hilb}}
\def\Curvesstack{\mathcal{C}\!\mathit{urves}}
\def\Polarizedstack{\mathcal{P}\!\mathit{olarized}}
\def\Complexesstack{\mathcal{C}\!\mathit{omplexes}}
% \Pic is the operator that assigns to X its picard group, usage \Pic(X)
% \Picardstack_{X/B} denotes the Picard stack of X over B
% \Picardfunctor_{X/B} denotes the Picard functor of X over B
\def\Pic{\mathop{\mathrm{Pic}}\nolimits}
\def\Picardstack{\mathcal{P}\!\mathit{ic}}
\def\Picardfunctor{\mathrm{Pic}}
\def\Deformationcategory{\mathcal{D}\!\mathit{ef}}
\usepackage{cleveref}
\crefname{originalreference}{original reference}{original references}
\Crefname{originalreference}{Original reference}{Original references}
\makeatletter
\newlabel{lemma-derived-quasi-coherent-small-etale-site}{{071Q}{source}{}{sourcepointer-1}{}}
\newlabel{lemma-derived-quasi-coherent-small-etale-site@cref}{{[originalreference][0][]071Q}{[1][0][]source}{}{}{}}
\makeatother
\crefname{originalreference}{original reference}{original references}
\Crefname{originalreference}{Original reference}{Original references}
\makeatletter
\newlabel{lemma-quasi-coherence-pushforward-direct-sums}{{08DZ}{source}{}{sourcepointer-2}{}}
\newlabel{lemma-quasi-coherence-pushforward-direct-sums@cref}{{[originalreference][0][]08DZ}{[1][0][]source}{}{}{}}
\makeatother
\crefname{originalremark}{original remark}{original remarks}
\Crefname{originalremark}{Original remark}{Original remarks}
\makeatletter
\newlabel{more-algebra-remark-P-resolution}{{09PB}{source}{}{sourcepointer-3}{}}
\newlabel{more-algebra-remark-P-resolution@cref}{{[originalremark][0][]09PB}{[1][0][]source}{}{}{}}
\makeatother
\crefname{originallemma}{original lemma}{original lemmas}
\Crefname{originallemma}{Original lemma}{Original lemmas}
\makeatletter
\newlabel{perfect-lemma-affine-compare-bounded}{{06Z0}{source}{}{sourcepointer-4}{}}
\newlabel{perfect-lemma-affine-compare-bounded@cref}{{[originallemma][0][]06Z0}{[1][0][]source}{}{}{}}
\makeatother
\crefname{originallemma}{original lemma}{original lemmas}
\Crefname{originallemma}{Original lemma}{Original lemmas}
\makeatletter
\newlabel{sites-cohomology-lemma-adjoint}{{07A6}{source}{}{sourcepointer-5}{}}
\newlabel{sites-cohomology-lemma-adjoint@cref}{{[originallemma][0][]07A6}{[1][0][]source}{}{}{}}
\makeatother
\crefname{originallemma}{original lemma}{original lemmas}
\Crefname{originallemma}{Original lemma}{Original lemmas}
\makeatletter
\newlabel{sites-cohomology-lemma-pullback-tensor-product}{{07A4}{source}{}{sourcepointer-6}{}}
\newlabel{sites-cohomology-lemma-pullback-tensor-product@cref}{{[originallemma][0][]07A4}{[1][0][]source}{}{}{}}
\makeatother
\title{The projection formula for derived pushforward of quasi-coherent complexes}
\author{The Stacks Project Authors}
\date{stacks-project:spaces-perfect.tex:local-snapshot-20261012; GNU FDL 1.2 or later}
\begin{document}
\maketitle
\noindent\textit{Editorial note.} All delivered mathematical text is the author's own, copied from the pinned Stacks Project chapter file spaces-perfect.tex; no statement or proof step was rewritten, reordered or removed.\par\medskip
\noindent\textit{Editorial note.} The counted proof is the author's complete argument: the proof is complete: the unit map, the sum and triangle compatibilities and the affine devissage to shifts of the structure sheaf\par\medskip
\noindent\textit{Editorial note.} The label lemma-cohomology-base-change names the counted statement, and the proof environment immediately following it in the source is delivered verbatim as its proof.\par\medskip
\noindent\textit{Editorial note.} The counted statement is carried by the author's own lemma environment in the chapter source.\par\medskip
\noindent\textit{Allowed standard prerequisites (editorial).}\begin{itemize}
\item derived pushforward along a quasi-compact quasi-separated morphism of algebraic spaces
\item derived tensor and derived pullback and their compatibility under f\textasciicircum{}*
\item commutation of Rf\_* with direct sums for complexes with quasi-coherent cohomology
\item triangulated functoriality of the square
\item the affine equivalence D(A) tilde-equivalent D\_QCoh(O\_Y) and the P-resolution remark reducing to shifts of the structure sheaf
\end{itemize}\medskip
\ifdefined\hypertarget\hypertarget{sourcepointer-1}{}\fi
\noindent\textit{Source reference.} Original reference 071Q: sources/stacks-spaces-perfect/pinned-original-source.tex, source lines 402--402. Prerequisite pointer invoked at proof line(s) 4929; it names a standard fact of the same chapter that the counted argument uses, and every step of the counted proof built on it is retained verbatim in the delivered body.\par\medskip
\ifdefined\hypertarget\hypertarget{sourcepointer-2}{}\fi
\noindent\textit{Source reference.} Original reference 08DZ: sources/stacks-spaces-perfect/pinned-original-source.tex, source lines 863--863. Prerequisite pointer invoked at proof line(s) 4919; it names a standard fact of the same chapter that the counted argument uses, and every step of the counted proof built on it is retained verbatim in the delivered body.\par\medskip
\ifdefined\hypertarget\hypertarget{sourcepointer-3}{}\fi
\noindent\textit{Source reference.} Original remark 09PB: sources/stacks-chow/tags.txt, source lines 12244--12244. Prerequisite pointer invoked at proof line(s) 4935; it names a standard fact of a sibling chapter that the counted argument uses, and every step of the counted proof built on it is retained verbatim in the delivered body.\par\medskip
\ifdefined\hypertarget\hypertarget{sourcepointer-4}{}\fi
\noindent\textit{Source reference.} Original lemma 06Z0: sources/stacks-chow/tags.txt, source lines 8904--8904. Prerequisite pointer invoked at proof line(s) 4931; it names a standard fact of a sibling chapter that the counted argument uses, and every step of the counted proof built on it is retained verbatim in the delivered body.\par\medskip
\ifdefined\hypertarget\hypertarget{sourcepointer-5}{}\fi
\noindent\textit{Source reference.} Original lemma 07A6: sources/stacks-chow/tags.txt, source lines 9296--9296. Prerequisite pointer invoked at proof line(s) 4908; it names a standard fact of a sibling chapter that the counted argument uses, and every step of the counted proof built on it is retained verbatim in the delivered body.\par\medskip
\ifdefined\hypertarget\hypertarget{sourcepointer-6}{}\fi
\noindent\textit{Source reference.} Original lemma 07A4: sources/stacks-chow/tags.txt, source lines 9294--9294. Prerequisite pointer invoked at proof line(s) 4907; it names a standard fact of a sibling chapter that the counted argument uses, and every step of the counted proof built on it is retained verbatim in the delivered body.\par\medskip
\begin{lemma}
\label{lemma-cohomology-base-change}
Let $S$ be a scheme. Let $f : X \to Y$ be a quasi-compact and quasi-separated
morphism of algebraic spaces over $S$. For $E$ in
$D_\QCoh(\mathcal{O}_X)$ and
$K$ in $D_\QCoh(\mathcal{O}_Y)$ we have
$$
Rf_*(E) \otimes_{\mathcal{O}_Y}^\mathbf{L} K =
Rf_*(E \otimes_{\mathcal{O}_X}^\mathbf{L} Lf^*K)
$$
\end{lemma}

\begin{proof}
Without any assumptions there is a map
$Rf_*(E) \otimes_{\mathcal{O}_Y}^\mathbf{L} K \to
Rf_*(E \otimes_{\mathcal{O}_X}^\mathbf{L} Lf^*K)$.
Namely, it is the adjoint to the canonical map
$$
Lf^*(Rf_*(E) \otimes_{\mathcal{O}_Y}^\mathbf{L} K) =
Lf^*(Rf_*(E)) \otimes_{\mathcal{O}_X}^\mathbf{L} Lf^*K
\longrightarrow
E \otimes_{\mathcal{O}_X}^\mathbf{L} Lf^*K
$$
coming from the map $Lf^*Rf_*E \to E$. See
Cohomology on Sites, Lemmas
\ref{sites-cohomology-lemma-pullback-tensor-product} and
\ref{sites-cohomology-lemma-adjoint}.
To check it is an isomorphism we may work \'etale locally on $Y$.
Hence we reduce to the case that $Y$ is an affine scheme.

\medskip\noindent
Suppose that $K = \bigoplus K_i$ is a direct
sum of some complexes $K_i \in D_\QCoh(\mathcal{O}_Y)$.
If the statement holds for each $K_i$, then it holds for $K$.
Namely, the functors $Lf^*$ and $\otimes^\mathbf{L}$ preserve
direct sums by construction and $Rf_*$ commutes with direct sums
(for complexes with quasi-coherent cohomology sheaves) by
Lemma \ref{lemma-quasi-coherence-pushforward-direct-sums}.
Moreover, suppose that $K \to L \to M \to K[1]$ is a distinguished
triangle in $D_\QCoh(Y)$. Then if the statement of the
lemma holds for two of $K, L, M$, then it holds for the third
(as the functors involved are exact functors of triangulated categories).

\medskip\noindent
Assume $Y$ affine, say $Y = \Spec(A)$. The functor
$\widetilde{\ } : D(A) \to D_\QCoh(\mathcal{O}_Y)$ is an equivalence
by
Lemma \ref{lemma-derived-quasi-coherent-small-etale-site} and
Derived Categories of Schemes,
Lemma \ref{perfect-lemma-affine-compare-bounded}.
Let $T$ be the property for $K \in D(A)$ that
the statement of the lemma holds for $\widetilde{K}$.
The discussion above and
More on Algebra, Remark \ref{more-algebra-remark-P-resolution}
shows that it suffices to prove $T$ holds for $A[k]$.
This finishes the proof, as the statement of the lemma
is clear for shifts of the structure sheaf.
\end{proof}

\end{document}
